\documentclass[11pt]{article}
\usepackage[margin=1in]{geometry}
\usepackage{amsmath,amssymb}
\usepackage{enumitem}
\title{MATH 6361 Mathematical Logic: Module Summary}
\author{Timothy Trujillo}
\date{}
\begin{document}
\maketitle

\section*{Course Architecture}
The course is organized as a module sequence. Each module combines guiding questions, learning outcomes, reading notes, video notes, homework, discussion, and quiz practice. The central movement is from syntax and semantics to proof systems, metatheory, quantification, set-theoretic structures, completeness, and incompleteness.

\section{Module 0: Atomic Sentences}
\textbf{Objectives.} Identify terms, predicates, function symbols, variables, constants, and atomic sentences. Translate between English and first-order atomic sentences.
\[
P(t_1,\ldots,t_n)
\]
is the basic form of an atomic sentence.

\section{Modules 1--3: Propositional Logic and Boolean Connectives}
\textbf{Objectives.} Classify arguments as valid or invalid and sound or unsound. Construct truth tables for \(\neg,\wedge,\vee,\to,\leftrightarrow\). Distinguish logical truth, tautology, consequence, and equivalence.

\section{Modules 4--5: Formal Proof Methods}
\textbf{Objectives.} Use introduction and elimination rules, subproofs, proof by contradiction, proof by cases, conditionals, and biconditionals in Fitch-style derivations.

\section{Module 6: Soundness and Completeness}
\textbf{Objectives.} State and interpret:
\[
\Gamma \vdash \varphi \Rightarrow \Gamma \models \varphi
\]
and
\[
\Gamma \models \varphi \Rightarrow \Gamma \vdash \varphi.
\]
Students learn to treat proof systems as mathematical objects.

\section{Module 9: Quantification}
\textbf{Objectives.} Recognize well-formed formulas, sentences, free variables, bound variables, and truth of quantified statements in structures.
\[
\forall x\,P(x), \qquad \exists x\,P(x).
\]

\section{Modules 10--11: Proof Methods for Quantifiers}
\textbf{Objectives.} Translate statements with multiple quantifiers, analyze quantifier order, and construct formal proofs using universal and existential introduction and elimination.
\[
\forall x\exists y\,R(x,y) \quad \text{need not mean} \quad \exists y\forall x\,R(x,y).
\]

\section{Module 12: Axiomatic Set Theory and Structures}
\textbf{Objectives.} Apply extension, separation, power set, union, intersection, relation, and function concepts. Use set theory as a language for mathematical structures.

\section{Modules 13--14: Completeness and Incompleteness}
\textbf{Objectives.} State soundness, completeness, and incompleteness theorems. Explain the tension between formal proof and mathematical truth in sufficiently expressive systems.

\end{document}
